%1/2in left/right margins, 1in top, 1/2 bottom
\documentclass[12pt]{article}
\hbadness=10000
\tolerance=10000
\textheight 9.7in
\textwidth 7.0in
\oddsidemargin -.5in
\topmargin -.5in
\headsep 0in
\headheight 0in
%\pagestyle{empty}

\newcommand{\by}{\mbox{\boldmath$y$}}
\newcommand{\bX}{\mbox{\boldmath$X$}}
\newcommand{\bbeta}{\mbox{\boldmath$\beta$}}

\begin{document}
\hrule
\medskip
\centerline{\textbf{STATISTICS 801 - Section 250}}\hfill \textbf{Fall 2016}
\centerline{\textbf{Mini Project 1}}
\smallskip
\hrule
\begin{enumerate}
\item Find some numerical data that interest you. The source can be your research, Internet, publication, etc. Document where you found the data.
\item Provide the data - either print, or print subset of the data. If you provide subset of the data, please provide a subset that is representative of the data, and helps understand the data in general. 
\item Explain why and how the data were collected.
\item What is the population in your example? What is the sample?
\item List the variables in your example.
\item Using at least two graphical methods summarize your data. You don't have to use the same graphs as presented in the lecture notes.
\item Explain why the graphs you chose in part 6. describe the data well.
\item Provide a graph that is a misrepresentation of your data. Explain why you think that.
\end{enumerate}

\textbf{You can use any software for this project! Do NOT draw graphs by hand!}
\end{document}

